\documentclass{article}
\usepackage[utf8]{inputenc}
\usepackage[vietnamese,english]{babel}
\usepackage{times}
\usepackage{amsmath,amssymb}
\usepackage{tabularx}
\usepackage{array}
\usepackage{booktabs}
\usepackage{graphicx}
\usepackage{xcolor}
\usepackage[
  paperwidth=29.7cm,
  paperheight=16.706cm,
  top=1.0cm,
  bottom=1.0cm,
  left=1.2cm,
  right=1.2cm,
  footskip=0pt,
  headheight=0pt,
  headsep=0pt
]{geometry}

\pagestyle{empty}
\setlength{\parindent}{0pt}
\setlength{\parskip}{0pt}

% Định nghĩa màu sắc thương hiệu
\definecolor{brandGreen}{RGB}{46, 125, 50}
\definecolor{brandRed}{RGB}{211, 47, 47}
\definecolor{brandBlue}{RGB}{21, 101, 192}

\begin{document}
\selectlanguage{vietnamese}

% ==============================================================================
% SLIDE 1: TRANG TIÊU ĐỀ
% ==============================================================================
\begin{center}
\vspace*{0.6cm}
{\fontsize{24}{30}\selectfont \textbf{TỔ KHOA HỌC TỰ NHIÊN}}\\[0.6cm]
{\fontsize{34}{42}\selectfont \textbf{\color{brandGreen}HƯỚNG DẪN ÔN TẬP GIỮA HỌC KỲ I}}\\[0.4cm]
{\fontsize{28}{36}\selectfont \textbf{MÔN: KHOA HỌC TỰ NHIÊN 8 -- NĂM HỌC 2026 - 2027}}\\[1.0cm]
{\fontsize{24}{32}\selectfont \textit{Phân môn: Hóa học (Chất và biến đổi) -- Sinh học (Vật sống) -- Vật lí (Năng lượng và biến đổi)}}\\[1.2cm]
{\fontsize{24}{30}\selectfont \textbf{\color{brandRed}Thầy TRẦN VĂN THẠNH -- 0777.470.803}}\\[0.3cm]
{\fontsize{20}{26}\selectfont CS1: 6/15 Nguyễn Hoàng, P. Kim Long, TP Huế \quad | \quad CS2: 24 Đặng Thái Thân, TP Huế}
\end{center}
\newpage

% ==============================================================================
% PHẦN I: HÓA HỌC - TRẮC NGHIỆM
% ==============================================================================

% SLIDE 2: HÓA CÂU 1 & CÂU 2
\fontsize{26}{34}\selectfont
\textbf{Câu 1.} Chất được tạo thành sau phản ứng hóa học là?
\vspace{6pt}

\begin{tabularx}{\linewidth}{@{}X X X X@{}}
\textbf{A.} Chất phản ứng. & \textbf{B.} Chất lỏng. & \textbf{C.} Chất sản phẩm. & \textbf{D.} Chất khí.
\end{tabularx}

\vspace*{\fill}

\textbf{Câu 2.} Phản ứng tỏa nhiệt là:
\vspace{6pt}

\textbf{A.} Phản ứng có nhiệt độ lớn hơn môi trường xung quanh.\\
\textbf{B.} Phản ứng có nhiệt độ nhỏ hơn môi trường xung quanh.\\
\textbf{C.} Phản ứng có nhiệt độ bằng môi trường xung quanh.\\
\textbf{D.} Phản ứng không có sự thay đổi nhiệt độ.
\vspace*{\fill}
\newpage

% SLIDE 3: HÓA CÂU 3 & CÂU 4
\fontsize{26}{34}\selectfont
\textbf{Câu 3.} Hòa tan đường vào nước là:
\vspace{6pt}

\begin{tabularx}{\linewidth}{@{}X X@{}}
\textbf{A.} Phản ứng hóa học. & \textbf{B.} Phản ứng tỏa nhiệt.\\[8pt]
\textbf{C.} Phản ứng thu nhiệt. & \textbf{D.} Sự biến đổi vật lí.
\end{tabularx}

\vspace*{\fill}

\textbf{Câu 4.} Chất mới được tạo ra từ phản ứng hóa học so với chất cũ sẽ như thế nào?
\vspace{6pt}

\begin{tabularx}{\linewidth}{@{}X X@{}}
\textbf{A.} Có tính chất mới, khác biệt chất ban đầu. & \textbf{B.} Giống hệt chất ban đầu.\\[8pt]
\textbf{C.} A và B đều đúng. & \textbf{D.} A và B đều sai.
\end{tabularx}
\vspace*{\fill}
\newpage

% SLIDE 4: HÓA CÂU 5
\fontsize{26}{34}\selectfont
\textbf{Câu 5.} Than (thành phần chính là carbon) cháy trong không khí tạo thành khí carbon dioxide. Trong quá trình phản ứng, lượng chất nào tăng dần?
\vspace{16pt}

\begin{tabularx}{\linewidth}{@{}X X@{}}
\textbf{A.} Carbon dioxide tăng dần. & \textbf{B.} Oxygen tăng dần.\\[12pt]
\textbf{C.} Carbon tăng dần. & \textbf{D.} Tất cả đều tăng.
\end{tabularx}
\vspace*{\fill}
\newpage

% SLIDE 5: HÓA ĐÚNG/SAI CÂU 1
\fontsize{24}{32}\selectfont
\textbf{Mục 2.2.} Xác định các nhận định sau về Mol và chất khí là Đúng (Đ) hay Sai (S):
\vspace{12pt}

a) Mol là lượng chất có chứa $N_A$ nguyên tử hoặc phân tử của chất đó. \hfill $[\quad]$\\[10pt]
b) Khối lượng mol của một chất là khối lượng tính bằng gam của $N_A$ nguyên tử hoặc phân tử chất đó. \hfill $[\quad]$\\[10pt]
c) Thể tích của 1 mol bất kỳ chất khí nào ở điều kiện chuẩn cũng luôn bằng $22,4\text{ lít}$. \hfill $[\quad]$\\[10pt]
d) Ở cùng điều kiện nhiệt độ và áp suất, 1 mol của mọi chất khí đều chiếm những thể tích bằng nhau. \hfill $[\quad]$\\[10pt]
e) Trong $0,12\text{ mol}$ phân tử $\text{NH}_3$ có chứa $0,24\text{ mol}$ nguyên tử N. \hfill $[\quad]$
\vspace*{\fill}
\newpage

% SLIDE 6: HÓA TRẢ LỜI NGẮN (Ý a & b)
\fontsize{26}{34}\selectfont
\textbf{a)} Cho sơ đồ: $\text{Fe} + \text{S} \xrightarrow{t^\circ} \text{FeS}$. Nếu dùng $7\text{g}$ sắt tác dụng vừa đủ với $4\text{g}$ lưu huỳnh thì khối lượng muối sắt(II) sulfide ($\text{FeS}$) thu được là bao nhiêu gam?\\
\textbf{Trả lời:} $\dots\dots\dots\dots\dots\dots\dots\dots\dots\dots$

\vspace*{\fill}

\textbf{b)} Nung nóng hoàn toàn $12\text{g}$ đá vôi $\text{CaCO}_3$ thu được $6,72\text{g}$ khí $\text{CO}_2$ và một lượng chất rắn $\text{CaO}$. Khối lượng của $\text{CaO}$ thu được là bao nhiêu gam?\\
\textbf{Trả lời:} $\dots\dots\dots\dots\dots\dots\dots\dots\dots\dots$
\vspace*{\fill}
\newpage

% SLIDE 7: HÓA TRẢ LỜI NGẮN (Ý c)
\fontsize{26}{34}\selectfont
\textbf{c)} Đốt cháy hoàn toàn $48\text{g}$ lưu huỳnh ($\text{S}$) trong khí oxygen thu được $96\text{g}$ khí sulfur dioxide ($\text{SO}_2$). Khối lượng khí oxygen ($\text{O}_2$) đã tham gia phản ứng là bao nhiêu gam?
\vspace{16pt}

\textbf{Trả lời:} $\dots\dots\dots\dots\dots\dots\dots\dots\dots\dots$
\vspace*{\fill}
\newpage

% SLIDE 8: HÓA TỰ LUẬN CÂU 1
\fontsize{26}{34}\selectfont
\textbf{Câu 1 (Tự luận Hóa).} Xác định độ tan của muối $\text{Na}_2\text{CO}_3$ trong nước ở $18^\circ\text{C}$. Biết rằng ở nhiệt độ này khi hòa tan hết $53\text{g } \text{Na}_2\text{CO}_3$ trong $250\text{g}$ nước thì được dung dịch bão hòa.
\vspace*{\fill}
\newpage

% SLIDE 9: HÓA TỰ LUẬN CÂU 2
\fontsize{25}{33}\selectfont
\textbf{Câu 2 (Tự luận Hóa).} Hòa tan $8\text{ gam}$ đường vào nước được dung dịch có nồng độ $5\%$. Hãy tính:\\
a. Khối lượng dung dịch đường pha chế được.\\
b. Khối lượng nước cần dùng cho sự pha chế.\\
c. Cần thêm vào bao nhiêu gam nước để được dung dịch có nồng độ $4\%$.\\
d. Cần thêm vào bao nhiêu gam đường để được dung dịch có nồng độ $6\%$.
\vspace*{\fill}
\newpage

% SLIDE 10: HÓA TỰ LUẬN CÂU 3
\fontsize{26}{34}\selectfont
\textbf{Câu 3 (Tự luận Hóa).} Tính số mol và số gam chất tan có trong $250\text{ mL}$ dung dịch $\text{CaCl}_2$ $0,3\text{M}$. (Cho biết $\text{Cl} = 35,5; \text{Ca} = 40$).
\vspace*{\fill}
\newpage

% SLIDE 11: HÓA TỰ LUẬN CÂU 4
\fontsize{25}{33}\selectfont
\textbf{Câu 4 (Tự luận Hóa).} Trong $800\text{ mL}$ của một dung dịch có chứa $8\text{g } \text{NaOH}$.\\
a. Hãy tính nồng độ mol của dung dịch này.\\
b. Phải thêm bao nhiêu $\text{mL}$ nước vào $200\text{ mL}$ dung dịch này để được dung dịch $\text{NaOH}$ $0,1\text{M}$.
\vspace*{\fill}
\newpage

% SLIDE 12: HÓA TỰ LUẬN CÂU 5
\fontsize{25}{33}\selectfont
\textbf{Câu 5 (Tự luận Hóa).} Chọn hệ số và CTHH thích hợp đặt vào dấu hỏi:\\
\begin{tabularx}{\linewidth}{@{}X X@{}}
a. $? + ? \xrightarrow{t^\circ} 2\text{FeCl}_3$ & b. $\text{Fe}_x\text{O}_y + ?\text{H}_2 \xrightarrow{t^\circ} ?\text{Fe} + ?\text{H}_2\text{O}$\\[8pt]
c. $?\text{Al} + ? \xrightarrow{t^\circ} 2\text{Al}_2\text{O}_3$ & d. $?\text{Ca(OH)}_2 + ?\text{P}_2\text{O}_5 \to ?\text{Ca}_3(\text{PO}_4)_2 + ?\text{H}_2\text{O}$\\[8pt]
e. $\text{C}_n\text{H}_{2n+2} + ?\text{O}_2 \xrightarrow{t^\circ} ?\text{CO}_2 + ?\text{H}_2\text{O}$ & g. $?\text{Al} + ?\text{Cu}_x\text{O}_y \xrightarrow{t^\circ} ?\text{Al}_2\text{O}_3 + ?\text{Cu}$
\end{tabularx}
\vspace*{\fill}
\newpage

% SLIDE 13: HÓA TỰ LUẬN CÂU 6
\fontsize{26}{34}\selectfont
\textbf{Câu 6 (Tự luận Hóa).} Tính số mol và số gam chất tan có trong $200\text{ mL}$ dung dịch $\text{KCl}$ $0,25\text{M}$. (Cho biết $\text{Cl} = 35,5, \text{K} = 39$).
\vspace*{\fill}
\newpage

% SLIDE 14: HÓA TỰ LUẬN CÂU 7
\fontsize{25}{33}\selectfont
\textbf{Câu 7 (Tự luận Hóa).} Trong $400\text{ mL}$ của một dung dịch có chứa $16\text{g } \text{CuSO}_4$.\\
a. Hãy tính nồng độ mol của dung dịch này.\\
b. Phải thêm bao nhiêu $\text{mL}$ nước vào $200\text{ mL}$ dung dịch này để được dung dịch $\text{CuSO}_4$ $0,2\text{M}$.
\vspace*{\fill}
\newpage

% SLIDE 15: HÓA TỰ LUẬN CÂU 8
\fontsize{25}{33}\selectfont
\textbf{Câu 8 (Tự luận Hóa).} Hãy tính toán và trình bày cách pha chế $50\text{ gam}$ dung dịch $\text{NaCl}$ có nồng độ $2\%$. (Dụng cụ và hoá chất coi như có đủ).
\vspace*{\fill}
\newpage

% ==============================================================================
% PHẦN III: VẬT LÍ - KHỐI LƯỢNG RIÊNG VÀ ÁP SUẤT
% ==============================================================================

% SLIDE 16: VẬT LÍ CÂU 1 & CÂU 2
\fontsize{26}{34}\selectfont
\textbf{Câu 1.} Đơn vị đo của khối lượng riêng trong hệ SI là:
\vspace{6pt}

\begin{tabularx}{\linewidth}{@{}X X X X@{}}
\textbf{A.} $\text{kg/m}^3$. & \textbf{B.} $\text{N}$. & \textbf{C.} $\text{kg}$. & \textbf{D.} $\text{N/m}^2$.
\end{tabularx}

\vspace*{\fill}

\textbf{Câu 2.} Muốn đo khối lượng riêng của viên bi người ta dùng những dụng cụ gì?
\vspace{6pt}

\textbf{A.} Chỉ cần dùng một cái cân.\\
\textbf{B.} Chỉ cần dùng một lực kế.\\
\textbf{C.} Cần dùng một cái cân và bình chia độ.\\
\textbf{D.} Chỉ cần dùng một bình chia độ.
\vspace*{\fill}
\newpage

% SLIDE 17: VẬT LÍ CÂU 3 & CÂU 4
\fontsize{25}{33}\selectfont
\textbf{Câu 3.} Đơn vị của áp suất là:
\vspace{6pt}

\begin{tabularx}{\linewidth}{@{}X X X X@{}}
\textbf{A.} $\text{Pa}$. & \textbf{B.} $\text{N}$. & \textbf{C.} $\text{m}$. & \textbf{D.} $\text{N/m}^3$.
\end{tabularx}

\vspace*{\fill}

\textbf{Câu 4.} Để xác định khối lượng riêng của các viên sỏi, ba bạn Sử, Sen, Anh đưa ra ý kiến. Theo em, ý kiến của bạn nào đúng?
\vspace{6pt}

\begin{tabularx}{\linewidth}{@{}X X@{}}
\textbf{A.} Sử đúng (chỉ cần cân). & \textbf{B.} Sen đúng (chỉ cần bình chia độ).\\[8pt]
\textbf{C.} Anh đúng (cần cả cân và bình chia độ). & \textbf{D.} Cả ba bạn cùng sai.
\end{tabularx}
\vspace*{\fill}
\newpage

% SLIDE 18: VẬT LÍ CÂU 5 & CÂU 6
\fontsize{25}{33}\selectfont
\textbf{Câu 5.} Khối lượng riêng của sắt là $7800\text{ kg/m}^3$. Vậy, $1\text{kg}$ sắt sẽ có thể tích vào khoảng
\vspace{6pt}

\begin{tabularx}{\linewidth}{@{}X X X X@{}}
\textbf{A.} $12,8\text{ cm}^3$. & \textbf{B.} $128\text{ cm}^3$. & \textbf{C.} $1280\text{ cm}^3$. & \textbf{D.} $12800\text{ cm}^3$.
\end{tabularx}

\vspace*{\fill}

\textbf{Câu 6.} Phát biểu nào sau đây về khối lượng riêng là đúng?
\vspace{6pt}

\textbf{A.} Khối lượng riêng của một chất là khối lượng của một đơn vị thể tích chất đó.\\
\textbf{B.} Nói khối lượng riêng của sắt là $7800\text{ kg/m}^3$ có nghĩa là $1\text{ cm}^3$ sắt có khối lượng $7800\text{ kg}$.\\
\textbf{C.} Công thức tính khối lượng riêng là $D = m \cdot V$.\\
\textbf{D.} Khối lượng riêng bằng trọng lượng riêng.
\vspace*{\fill}
\newpage

% SLIDE 19: VẬT LÍ CÂU 7 & CÂU 8
\fontsize{25}{33}\selectfont
\textbf{Câu 7.} Muốn tăng áp suất tác dụng lên mặt bị ép, ta có thể:
\vspace{6pt}

\textbf{A.} Giảm áp lực hoặc tăng diện tích mặt bị ép.\\
\textbf{B.} Tăng độ lớn áp lực hoặc giảm diện tích mặt bị ép.\\
\textbf{C.} Giữ nguyên áp lực và tăng diện tích mặt bị ép.\\
\textbf{D.} Giảm đồng thời áp lực và diện tích mặt bị ép.

\vspace*{\fill}

\textbf{Câu 8.} \textbf{Chỉ ra kết luận sai} trong các kết luận sau:
\vspace{6pt}

\textbf{A.} Áp lực là lực ép có phương vuông góc với mặt bị ép.\\
\textbf{B.} Đơn vị của áp suất là $\text{N/m}^2$.\\
\textbf{C.} Áp suất là độ lớn của áp lực trên một diện tích bị ép.\\
\textbf{D.} Đơn vị của áp lực là đơn vị của lực.
\vspace*{\fill}
\newpage

% SLIDE 20: VẬT LÍ TỰ LUẬN CÂU 9
\fontsize{25}{33}\selectfont
\textbf{Câu 9 (Vật lí).} Đặt một bao gạo $25\text{kg}$ lên một ghế 4 chân có khối lượng $4\text{kg}$. Diện tích tiếp xúc với mặt đất của mỗi chân ghế là $8\text{ cm}^2$. Áp suất mà gạo và ghế tác dụng lên mặt đất là bao nhiêu?
\vspace*{\fill}
\newpage

% SLIDE 21: VẬT LÍ TỰ LUẬN CÂU 10
\fontsize{25}{33}\selectfont
\textbf{Câu 10 (Vật lí).} Một vật hình lập phương đặt trên mặt bàn nằm ngang, tác dụng lên mặt bàn một áp suất $p = 3600\text{ N/m}^2$. Biết khối lượng của vật là $14,4\text{ kg}$. Hỏi độ dài một cạnh của khối lập phương ấy là bao nhiêu?
\vspace*{\fill}
\newpage

% SLIDE 22: VẬT LÍ TỰ LUẬN CÂU 11
\fontsize{25}{33}\selectfont
\textbf{Câu 11 (Vật lí).} Một cái dầm sắt có thể tích là $60\text{ dm}^3$, biết khối lượng riêng của sắt là $7,8\text{ g/cm}^3$. Tính khối lượng và trọng lượng của dầm sắt này.
\vspace*{\fill}
\newpage

% SLIDE 23: VẬT LÍ TỰ LUẬN CÂU 12
\fontsize{26}{34}\selectfont
\textbf{Câu 12 (Vật lí).} Một đồng xu có khối lượng $0,9\text{ g}$, được làm từ hợp kim có khối lượng riêng là $5\text{ g/cm}^3$. Tính thể tích của đồng xu.
\vspace*{\fill}
\newpage

% SLIDE 24: VẬT LÍ TỰ LUẬN CÂU 13
\fontsize{25}{33}\selectfont
\textbf{Câu 13 (Vật lí).} Một khối sắt hình hộp chữ nhật có kích thước $4\text{ cm}, 5\text{ cm}, 7\text{ cm}$ và khối lượng $700\text{ g}$. Tính trọng lượng và khối lượng riêng của khối sắt đó.
\vspace*{\fill}
\newpage

% SLIDE 25: VẬT LÍ TỰ LUẬN CÂU 14
\fontsize{26}{34}\selectfont
\textbf{Câu 14 (Vật lí).} Khối lượng riêng của dầu ăn vào khoảng $800\text{ kg/m}^3$. Do đó, $2\text{ lít}$ dầu ăn sẽ có trọng lượng khoảng bao nhiêu?
\vspace*{\fill}
\newpage

% ==============================================================================
% SLIDE 26: BẢNG TỔNG HỢP ĐÁP ÁN TOÀN BỘ ĐỀ CƯƠNG
% ==============================================================================
\begin{center}
{\fontsize{24}{30}\selectfont \textbf{\color{brandGreen}BẢNG TỔNG HỢP ĐÁP ÁN ĐỀ CƯƠNG ÔN TẬP KHTN 8}}\\[0.4cm]
\end{center}

\fontsize{18}{24}\selectfont
\textbf{1. Trắc nghiệm nhiều lựa chọn:}\\
\begin{tabularx}{\linewidth}{|c|*{5}{>{\centering\arraybackslash}X|}|c|*{8}{>{\centering\arraybackslash}X|}}
\hline
\textbf{Hóa} & \textbf{1} & \textbf{2} & \textbf{3} & \textbf{4} & \textbf{5} & \textbf{Lí} & \textbf{1} & \textbf{2} & \textbf{3} & \textbf{4} & \textbf{5} & \textbf{6} & \textbf{7} & \textbf{8}\\
\hline
\textbf{Đ/A} & \textbf{\color{brandRed}C} & \textbf{\color{brandRed}A} & \textbf{\color{brandRed}D} & \textbf{\color{brandRed}A} & \textbf{\color{brandRed}A} & \textbf{Đ/A} & \textbf{\color{brandRed}A} & \textbf{\color{brandRed}C} & \textbf{\color{brandRed}A} & \textbf{\color{brandRed}C} & \textbf{\color{brandRed}B} & \textbf{\color{brandRed}A} & \textbf{\color{brandRed}B} & \textbf{\color{brandRed}C}\\
\hline
\end{tabularx}
\vspace{8pt}

\textbf{2. Trắc nghiệm Đúng / Sai Hóa học (Mục 2.2):}\\
$\bullet$ a - \textbf{Đúng}, b - \textbf{Đúng}, c - \textbf{Sai}, d - \textbf{Đúng}, e - \textbf{Sai}.
\vspace{6pt}

\textbf{3. Trắc nghiệm Trả lời ngắn Hóa học (Mục 2.3):}\\
$\bullet$ a) \textbf{11 gam} \quad | \quad b) \textbf{5,28 gam} \quad | \quad c) \textbf{48 gam}.
\vspace{6pt}

\textbf{4. Đáp số bài tập tính toán Vật lí (Câu 9 -- 14):}\\
$\bullet$ \textbf{Câu 9:} $p = \mathbf{90\,625\text{ Pa}}$ \quad | \quad \textbf{Câu 10:} $a = \mathbf{0,2\text{ m}}\ (20\text{ cm})$.\\
$\bullet$ \textbf{Câu 11:} $m = \mathbf{468\text{ kg}},\ P = \mathbf{4680\text{ N}}$ \quad | \quad \textbf{Câu 12:} $V = \mathbf{0,18\text{ cm}^3}$.\\
$\bullet$ \textbf{Câu 13:} $P = \mathbf{7\text{ N}},\ D = \mathbf{5\text{ g/cm}^3}\ (5000\text{ kg/m}^3)$ \quad | \quad \textbf{Câu 14:} $P = \mathbf{16\text{ N}}$.

\end{document}
